\section{Experimental Setup}
\label{app:setup}

\subsection{Metric definitions and uncertainty}

\paragraph{Task metrics.}
For Crafter, let $r_a$ denote the fraction of evaluated worlds
in which achievement $a$ is attained. We report arithmetic
achievement progress,
\begin{equation}
P_{\mathrm{Crafter}}
=\frac{100}{22}\sum_{a=1}^{22}r_a.
\end{equation}
Peak Progress is the highest progress observed across evolution
rounds. The geometric achievement score is a distinct metric,
\begin{equation}
S_{\mathrm{Crafter}}
=\exp\left(
\frac{1}{22}\sum_{a=1}^{22}\log(1+100r_a)
\right)-1.
\label{eq:crafter}
\end{equation}
These statistics are computed from achievement rates across
worlds; averaging per-world geometric scores yields a different
quantity.

ALFWorld and WebShop use binary task success rate. WebShop's
fractional reward, when reported, is listed separately.
HotpotQA and WebQuestions use normalized answer exact match
(EM) and token F1. Local FinQA evaluations use scale-normalized
numeric-answer accuracy with a 1\% relative tolerance, which
is distinct from program execution accuracy.

\paragraph{Uncertainty.}
We distinguish standard deviations across repeated runs from
bootstrap uncertainty over evaluation instances. The local
paired analyses use checkpoints trained with seed 42; their
instance-level uncertainty does not measure variation across
independently trained models. Paired bootstrap samples use the
same resampled instances for all compared methods. For Crafter,
we resample worlds with their complete achievement vectors and
recompute the aggregate statistic in each draw.
Appendix~\ref{app:protocol} specifies the resampling procedures,
including hierarchical uncertainty where applicable.
Published comparisons retain the uncertainty reported by
their sources.

\subsection{Memory, retention, and serving measurements}
\label{app:costs}

\paragraph{Memory accounting.}
We distinguish retrievable ETG records, non-retrievable
archived records, and separately stored training data,
including replay when instantiated. For consolidation
analyses, we report record counts and available byte
measurements before and after the update. Moving records
to an archive reduces the retrievable footprint but does
not reduce total retained storage. Capability--footprint
comparisons therefore identify which stores are included.
For hard-cap configurations, rejected admissions and forced
deletions are recorded separately from verified release.

\paragraph{Retention.}
In multi-round evaluations, let $t_s$ be the first round
in which skill $s$ is consolidated and $R_{t,s}$ its rollout
success at round $t$. We report final success $R_{T,s}$ and
the decline from its best observed post-consolidation success,
\begin{equation}
D_{T,s}
=\max_{t_s\leq t\leq T}R_{t,s}-R_{T,s}.
\end{equation}
These measurements use fresh evaluation episodes.
Action-probe scores and training loss provide diagnostics;
they do not replace rollout-based retention measurements.

\paragraph{Acquisition and serving cost.}
We account separately for student and teacher acquisition,
verification probes, and training compute. Serving
measurements report elapsed time, model calls, and input
and output tokens per task, with retrieval tokens identified
where available. Paired comparisons use the same evaluation
tasks and distinguish weights-only, retrieval-assisted,
and per-skill routed deployment. Token usage and wall-clock
latency are reported separately because lower token usage
does not necessarily imply faster execution.

\paragraph{Release decisions.}
A skill is approved for release only after satisfying its
verification criterion. Candidate-adapter performance or
evaluation with retained ETG alone does not authorize
release. Development-probe approval and the resulting
record movement are reported separately from held-out
task performance after release.